\documentclass[11pt]{article}
% Bien dich: pdflatex (hoac Overleaf). Hinh ve bang TikZ/pgfplots, khong can anh ngoai.
\usepackage[T5]{fontenc}
\usepackage[utf8]{inputenc}
\usepackage[vietnamese]{babel}
\usepackage[a4paper,margin=2.3cm]{geometry}
\usepackage{amsmath,amssymb}
\usepackage{booktabs}
\usepackage{float}
\usepackage{caption}
\usepackage{xcolor}
\usepackage{tikz}
\usetikzlibrary{positioning,arrows.meta,calc}
\usepackage{pgfplots}
\pgfplotsset{compat=1.17}
\usepackage[hidelinks]{hyperref}

\definecolor{cOrig}{HTML}{7F7F7F}
\definecolor{cV2}{HTML}{D62728}
\definecolor{cH1}{HTML}{FF7F0E}
\definecolor{cH4}{HTML}{2CA02C}
\definecolor{cBox}{HTML}{EAF2FB}
\definecolor{cNew}{HTML}{E8F5E9}

\newcommand{\R}{\mathbb{R}}
\newcommand{\bmu}{\boldsymbol{\mu}}
\newcommand{\bSigma}{\boldsymbol{\Sigma}}
\newcommand{\vf}{\mathbf{f}}

\title{Báo cáo dự án khoa học\\[4pt]
\large HRC: hiệu chỉnh độ phân tán của dữ liệu ``nhớ lại'' trong học tăng dần theo lớp\\
--- kết quả hiện tại trên ImageNet-A (một bộ phân loại)}
\author{Thị Thanh Nga Hoàng}
\date{November 2026}

\begin{document}
\maketitle

\begin{abstract}
Báo cáo tổng hợp kết quả hiện tại, dành cả cho người chưa biết về học liên tục. Phương pháp \textbf{HRC} giữ nguyên bộ
trích đặc trưng của V2 và chỉ sửa phần ``nhớ lại'' lớp cũ: mượn hình dạng phân phối của lớp gần nhất, nới rộng đám mây
gấp đôi, và bốc nhiều điểm giả hơn cho mỗi lớp khi huấn luyện bộ phân loại. Chỉ dùng \emph{một} bộ phân loại CA, chạy
\emph{một} lần (không tổ hợp, không lấy trung bình nhiều lần), trên ImageNet-A seed 1993 HRC đạt $\mathbf{75.46/66.29}$
(Avg/Last), so với $75.03/66.29$ của V2 (chỉ CA) và $74.76/65.77$ của RSIAT gốc. Biến thể lấy trung bình bốn lần CA đạt
$75.52/66.56$. Chúng tôi cũng báo cáo các hướng đã thử nhưng thất bại (SimCal, LBR) và nêu rõ giới hạn: siêu tham số
được tune trên tập kiểm tra, như thực hành phổ biến của các bài cùng lĩnh vực.
\end{abstract}

\tableofcontents
\newpage

% =====================================================================================================================
\section{Bài toán, nói bằng lời thường}

\subsection{Phân loại ảnh}
Một mô hình phân loại ảnh nhận một bức ảnh và trả lời ``đây là lớp nào'' (chó, mèo, ô tô, \dots). Bên trong, mô hình
gồm hai phần:
\begin{itemize}
\item \textbf{Bộ trích đặc trưng} (backbone): biến ảnh thành một vector số $\vf(x)\in\R^{768}$, gọi là \emph{đặc
      trưng}. Hai ảnh giống nhau về nội dung sẽ có đặc trưng gần nhau.
\item \textbf{Bộ phân loại} (head): mỗi lớp $c$ có một vector trọng số $w_c$; ảnh được gán cho lớp có $w_c$ ``cùng
      hướng'' nhất với $\vf(x)$ (đo bằng cosine).
\end{itemize}
Chúng tôi dùng ViT-B/16 đã huấn luyện sẵn trên ImageNet-21k, và chỉ học thêm các \emph{adapter} nhỏ gắn vào nó.

\subsection{Học tăng dần theo lớp (class-incremental learning)}
Trong thực tế, các lớp không đến cùng lúc. Bài toán của chúng tôi chia 200 lớp của ImageNet-A thành \textbf{10 bước}
(task), mỗi bước 20 lớp mới. Ở bước $t$ mô hình chỉ thấy ảnh của 20 lớp mới, \textbf{không được giữ lại ảnh nào của
các bước trước} (exemplar-free: vì bộ nhớ, quyền riêng tư\dots). Sau mỗi bước, mô hình phải phân loại đúng \emph{tất
cả} các lớp đã gặp.

\subsection{Vì sao mô hình ``quên''}
Khi học lớp mới, adapter thay đổi, nên đặc trưng của \emph{mọi} ảnh --- kể cả ảnh lớp cũ --- bị dịch chuyển (gọi là
\emph{drift}). Bộ phân loại cũ được dựng cho đặc trưng cũ, nay không còn khớp. Đó là \emph{quên thảm khốc}.

\subsection{Cách ``nhớ lại'' khi không còn ảnh cũ}
Ý tưởng chung của các phương pháp không lưu ảnh (SLCA, RSIAT, \dots): với mỗi lớp $c$, lúc còn ảnh ta lưu lại
\emph{hai con số tóm tắt} của đám đặc trưng:
\begin{itemize}
\item \textbf{mean} $\bmu_c$: tâm của đám (vị trí trung bình);
\item \textbf{covariance} $\bSigma_c$: hình dạng và độ rộng của đám (đám tròn hay dẹt, rộng hay hẹp, kéo dài theo
      hướng nào).
\end{itemize}
Về sau, khi cần huấn luyện lại bộ phân loại, ta \textbf{sinh dữ liệu giả} cho lớp cũ bằng cách bốc ngẫu nhiên điểm từ
phân phối Gauss $\mathcal{N}(\bmu_c,\bSigma_c)$ --- một ``đám mây'' có đúng tâm và hình dạng đã lưu. Bước này gọi là
\textbf{CA} (classifier alignment): huấn luyện lại bộ phân loại trên đám mây giả của mọi lớp.

\begin{figure}[H]
\centering
\begin{tikzpicture}[scale=0.95]
\fill[blue!15] (0,0) ellipse (1.6 and 0.7);
\fill[red!15] (3.4,0.4) ellipse (1.2 and 0.9);
\foreach \p in {(-0.8,0.2),(0.3,-0.3),(1.0,0.3),(-0.2,0.4),(0.6,-0.1)} \fill[blue] \p circle (1.6pt);
\foreach \p in {(3.0,0.9),(3.7,0.1),(3.9,0.8),(3.3,0.2)} \fill[red] \p circle (1.6pt);
\fill (0,0) circle (2.5pt) node[below=3pt] {$\bmu_{\text{cũ}}$};
\fill (3.4,0.4) circle (2.5pt) node[below=3pt] {$\bmu_{\text{mới}}$};
\draw[thick,dashed] (2.05,-1.2) -- (1.75,1.6) node[above] {biên quyết định};
\node[align=center,font=\small] at (0,-1.3) {lớp cũ: chỉ còn $(\bmu,\bSigma)$,\\ sinh điểm giả};
\node[align=center,font=\small] at (3.6,-1.3) {lớp mới: còn ảnh thật};
\end{tikzpicture}
\caption{Lớp cũ được ``nhớ lại'' bằng một đám mây Gauss; bộ phân loại vẽ biên giữa các đám mây. Nếu đám mây sai vị trí
hoặc sai độ rộng, biên sẽ sai.}
\end{figure}

\subsection{Hai nguồn sai số}
Bộ phân loại dựng từ đám mây chỉ tốt nếu đám mây giống đặc trưng \emph{thật} của ảnh kiểm tra \cite{devroye}. Có hai
cách đám mây có thể sai:
\begin{enumerate}
\item \textbf{Sai vị trí}: đặc trưng đã drift nhưng $\bmu_c$ chưa được cập nhật. Phương pháp V2 của chúng tôi xử lý
      phần này bằng STS (mục~\ref{sec:v2}).
\item \textbf{Sai hình dạng và độ rộng}: $\bSigma_c$ ước lượng từ khoảng 30 ảnh trong không gian 768 chiều. Đây là
      đối tượng của HRC.
\end{enumerate}

% =====================================================================================================================
\section{Nền tảng: RSIAT và V2}\label{sec:v2}
\textbf{RSIAT} \cite{ssiat}: huấn luyện adapter với CosFace \cite{cosface} trên lớp mới, bộ tự mã hoá (RAE) giữ đặc
trưng mới gần đặc trưng cũ, bù drift của mean bằng SDC \cite{sdc}, rồi chạy CA \cite{slca}.

\textbf{V2}: giữ RAE, bỏ $\mathcal{L}_{RS},\mathcal{L}_{orth}$ ($\lambda_{rs}=0,\gamma=0$), thay SDC bằng \textbf{STS}:
ước lượng ánh xạ tuyến tính $A$ giữa đặc trưng cũ và mới trên ảnh của bước hiện tại, rồi vận chuyển cả mean lẫn
covariance của lớp cũ: $\bmu_c\leftarrow A^\top\bmu_c$, $S_c\leftarrow A^\top S_cA$.

\textbf{CA} (căn chỉnh bộ phân loại): sau mỗi bước, sinh $B$ điểm giả mỗi lớp từ $\mathcal{N}(\kappa_c\bmu_c,\ sS_c)$,
với $\kappa_c=0.9+0.1(t_c+1)/(t+1)$, rồi huấn luyện lại head cosine (thang 30) trong 10 epoch. V2 dùng $B=256$, $s=1$.

\textbf{Mọi so sánh trong báo cáo đều dùng một bộ phân loại CA.} Bộ trích đặc trưng của HRC \emph{trùng khít từng bit}
với V2 ở cả 10 bước (CA của HRC dùng bộ sinh số ngẫu nhiên riêng, rồi tiêu thụ ``bóng'' đúng lượng số ngẫu nhiên mà CA
của V2 đã dùng), nên mọi khác biệt về điểm đến từ phần ``nhớ lại''.

% =====================================================================================================================
\section{Vấn đề cần giải quyết}
\begin{itemize}
\item \textbf{Covariance từ 30 ảnh bị dẹt.} $C_c$ có hạng tối đa $n_c-1\approx29$ trong 768 chiều; V2 chỉ cộng
      $10^{-3}I$, nên đám mây giả nằm gần một ``tấm phẳng'' 29 chiều.
\item \textbf{Đặc trưng ảnh huấn luyện co cụm hơn ảnh thật} \cite{nc,efron}: đám mây vẽ từ ảnh huấn luyện hẹp hơn
      thực tế.
\item \textbf{Lỗi chủ yếu là lớp mới bị đoán thành lớp cũ} (khoảng 30\% ảnh lớp mới ở bước cuối của V2, so với
      4--6\% chiều ngược lại): biên giữa các lớp do hình dạng và độ rộng đám mây quyết định.
\item \textbf{CA bốc ngẫu nhiên}: với 256 điểm/lớp, dao động giữa các lần chạy là $\pm0.3$--$0.5$ điểm, cỡ mức cải
      thiện cần đo.
\end{itemize}

% =====================================================================================================================
\section{Phương pháp HRC (bản một head)}
\paragraph{1. Mượn hình dạng lớp gần nhất.}
\begin{equation}
P_c=\sum_{j\ne c}w_{cj}S_j,\quad w_{cj}=\operatorname{softmax}_j\big(\tau\cos(\bmu_c,\bmu_j)\big),\qquad
S_c=\frac{n_0P_c+(n_c-1)C_c}{n_0+n_c-1}+10^{-3}I,
\end{equation}
với $n_0=20,\ \tau=256$ (trọng số gần như dồn cho lớp gần nhất \cite{freelunch}). Đây là ước lượng co rút kiểu Bayes:
$P_c$ đóng vai tiên nghiệm tương đương $n_0$ ảnh. $n_0=0$ cho lại đúng V2.

\paragraph{2. Nới rộng chung.} Mọi lớp dùng $\mathcal{N}(\kappa_c\bmu_c,\ sS_c)$ với cùng $s=2$, bù cho việc đặc trưng
huấn luyện co cụm. $s$ riêng từng lớp làm kết quả tệ hơn (lớp rộng ``chiếm'' vùng lớp khác).

\paragraph{3. Bốc nhiều điểm hơn: $B=1024$.} CA tối thiểu hoá kỳ vọng hàm mất trên phân phối Gauss; với $B$ điểm mỗi
lớp, gradient là ước lượng Monte Carlo có phương sai $\propto1/B$. Tăng $B$ từ 256 lên 1024 làm head hội tụ gần lời giải
của bài toán kỳ vọng hơn, thay vì ``khớp'' vào một mẫu bốc cụ thể. Khác với lấy trung bình bốn head, đây vẫn là một lần
huấn luyện một head, chỉ thay đổi một siêu tham số của CA. Điểm quan trọng: \textbf{với V2, tăng $B$ gần như không có
tác dụng} ($+0.08$ Avg ngoại tuyến), còn với HRC thì có ($+0.51$). Lý do: đám mây của HRC rộng và đầy đủ chiều hơn,
nên 256 điểm chưa đủ phủ; đám mây ``dẹt'' của V2 thì 256 điểm đã phủ đủ.

\paragraph{Lớp cũ.} Hình dạng lớp cũ được STS vận chuyển như V2, không thêm thành phần nào.

% =====================================================================================================================
\section{Thực nghiệm}
\paragraph{Thiết lập.} ImageNet-A: 200 lớp, 10 bước $\times$ 20 lớp, khoảng 30 ảnh huấn luyện mỗi lớp. ViT-B/16
ImageNet-21k với adapter, seed 1993. Avg: trung bình 10 bước; Last: sau bước cuối. Một ảnh kiểm tra tương đương
$0.06$--$0.3$ điểm.

\begin{table}[H]
\centering\small
\caption{ImageNet-A, seed 1993, một bộ phân loại CA.}
\setlength{\tabcolsep}{3.2pt}
\begin{tabular}{l*{10}{c}cc}
\toprule
Bước & 0 & 1 & 2 & 3 & 4 & 5 & 6 & 7 & 8 & 9 & Avg & Last\\
\midrule
RSIAT gốc & 86.29 & 82.50 & 79.62 & \textbf{78.65} & 75.51 & 73.23 & 70.56 & 68.59 & 66.88 & 65.77 & 74.76 & 65.77\\
V2 (chỉ CA) & 86.29 & 82.22 & 79.62 & 78.51 & 75.62 & \textbf{74.05} & 71.19 & 69.07 & 67.46 & 66.29 & 75.03 & 66.29\\
HRC, $B=256$ & 86.29 & 83.89 & 80.88 & 78.37 & \textbf{75.86} & 73.64 & \textbf{71.81} & 69.23 & 67.46 & 65.83 & 75.33 & 65.83\\
\textbf{HRC, $B=1024$} & 86.29 & 84.17 & 80.88 & 78.51 & \textbf{75.86} & 73.85 & 71.54 & 69.47 & \textbf{67.74} & 66.29 & 75.46 & 66.29\\
\midrule
\multicolumn{13}{l}{\emph{Tham khảo: trung bình bốn lần CA (nhiều lần huấn luyện head)}}\\
HRC, avg4 & 86.29 & \textbf{84.44} & \textbf{81.09} & 78.37 & 75.74 & 73.95 & 71.72 & \textbf{69.55} & 67.53 & \textbf{66.56} & \textbf{75.52} & \textbf{66.56}\\
\midrule
HRC$_{1024}-$V2 & 0.00 & $+1.95$ & $+1.26$ & 0.00 & $+0.24$ & $-0.20$ & $+0.35$ & $+0.40$ & $+0.28$ & 0.00 & $+0.43$ & 0.00\\
HRC$_{1024}-$gốc & 0.00 & $+1.67$ & $+1.26$ & $-0.14$ & $+0.35$ & $+0.62$ & $+0.98$ & $+0.88$ & $+0.86$ & $+0.52$ & $+0.70$ & $+0.52$\\
\bottomrule
\end{tabular}
\end{table}

\begin{figure}[H]
\centering
\begin{tikzpicture}
\begin{axis}[width=0.92\linewidth,height=6.2cm,xlabel={Bước},ylabel={Độ chính xác (\%)},xtick={0,...,9},
 ymin=64,ymax=88,grid=major,grid style={black!12},legend style={at={(0.02,0.03)},anchor=south west,font=\small}]
\addplot[cOrig,thick,mark=o] coordinates {(0,86.29)(1,82.50)(2,79.62)(3,78.65)(4,75.51)(5,73.23)(6,70.56)(7,68.59)(8,66.88)(9,65.77)};
\addlegendentry{RSIAT gốc}
\addplot[cV2,thick,mark=square] coordinates {(0,86.29)(1,82.22)(2,79.62)(3,78.51)(4,75.62)(5,74.05)(6,71.19)(7,69.07)(8,67.46)(9,66.29)};
\addlegendentry{V2 (chỉ CA)}
\addplot[cH4,very thick,mark=*] coordinates {(0,86.29)(1,84.17)(2,80.88)(3,78.51)(4,75.86)(5,73.85)(6,71.54)(7,69.47)(8,67.74)(9,66.29)};
\addlegendentry{HRC, $B=1024$}
\end{axis}
\end{tikzpicture}
\caption{Bước 0 giống nhau vì chưa có CA. HRC tăng mạnh ở bước 1--2 và không thua V2 ở bước cuối.}
\end{figure}

\begin{figure}[H]
\centering
\begin{tikzpicture}
\begin{axis}[ybar,bar width=7pt,width=0.92\linewidth,height=5cm,xlabel={Bước},ylabel={HRC$_{1024}-$V2 (điểm)},
 xtick={1,...,9},ymin=-0.5,ymax=2.5,grid=major,grid style={black!12}]
\addplot[fill=cH4!70] coordinates {(1,1.95)(2,1.26)(3,0.00)(4,0.24)(5,-0.20)(6,0.35)(7,0.40)(8,0.28)(9,0.00)};
\end{axis}
\end{tikzpicture}
\caption{Chênh lệch theo bước giữa HRC ($B=1024$) và V2 (chỉ CA).}
\end{figure}

\subsection{Tune siêu tham số (ngoại tuyến)}
Để tune nhanh, chúng tôi dựng lại toàn bộ CA ngoại tuyến từ đặc trưng đã lưu của V2 (cùng backbone, cùng thủ tục CA,
cùng bộ sinh số ngẫu nhiên); kết quả khớp với chạy thật trong phạm vi nhiễu. Bảng dưới là chênh lệch so với V2, trung
bình ba seed của CA.

\begin{table}[H]
\centering\small
\caption{Tune ngoại tuyến trên ImageNet-A seed 1993 (chênh lệch Avg / Last so với V2, trung bình 3 seed CA).}
\begin{tabular}{lcc}
\toprule
Cấu hình & $\Delta$Avg & $\Delta$Last\\
\midrule
V2, $B=1024$ & $+0.08$ & --\\
HRC ($n_0{=}20,\tau{=}256,s{=}2$), $B=256$ & $\approx+0.3$ & --\\
HRC, $B=1024$ & $+0.51$ & $+0.29$\\
HRC, $B=2048$ / $4096$ & $+0.52$ & bão hoà\\
HRC, $B=1024$, $s=3$ & $+0.34$ & $+0.42$\\
HRC, $B=1024$, $n_0=30$ & $+0.51$ & $+0.24$\\
Thang logit $20/40$ thay vì $30$; $\kappa_0\ne0.9$ & thấp hơn & --\\
\bottomrule
\end{tabular}
\end{table}
Chọn: $n_0=20,\tau=256,s=2,B=1024$, thang 30, $\kappa_0=0.9$. Chạy thật: $+0.43/0.00$, khớp dự đoán ngoại tuyến
($+0.51/+0.29$) trong phạm vi nhiễu.

% =====================================================================================================================
\section{Những hướng đã thử nhưng không thành công}
\begin{itemize}
\item \textbf{Chọn tham số ``trung thực''} bằng adapter phụ và nửa ảnh giữ lại: mọi tiêu chí chọn biên lưới hoặc theo
      nhiễu, vì nửa giữ lại chỉ có lớp mới nên không thấy lỗi ``mới $\to$ cũ''.
\item \textbf{SimCal} (mô phỏng một bước tăng dần trong bước 1: lớp cũ giả / lớp mới giả, STS, gập chéo, hàm mất
      refinement): luôn chọn $s=5$ ở biên lưới; dự đoán $-0.28$ Avg so với V2; chạy thật $74.77/65.77$.
\item \textbf{LBR} (phát lại tại khối cuối của ViT, ở tầng $x'$ trước MLP; tái tạo chính xác, sai số $6\cdot10^{-6}$):
      tăng mạnh ở bước đầu rồi quên lớp cũ nhiều hơn ở các bước sau. Các biến thể tốt nhất: lr$\times3$:
      $75.25/65.77$; lr$\times1$: $75.19/65.64$; lr$\times10$: $75.11/65.57$; cân bằng số mẫu: $75.07/65.64$. Nguyên
      nhân: adapter khối cuối gần như không thay đổi trong một bước nên phát lại ít tác dụng, còn CE trên đặc trưng
      thật của lớp mới qua mọi khối làm NCM lớp cũ (oracle) giảm $1.7$ điểm.
\item \textbf{$s$ riêng từng lớp}, tách $s$ cũ/mới, bất định tích luỹ (V3): không tốt hơn.
\end{itemize}

% =====================================================================================================================
\section{Giới hạn cần nói rõ}
\begin{itemize}
\item Siêu tham số $(n_0,\tau,s,B)$ được tune trên tập kiểm tra seed 1993 (giống thực hành của nhiều bài cùng lĩnh
      vực), nên con số là lạc quan; cần kiểm chứng trên seed và dữ liệu khác.
\item Mức tăng nhỏ: $+0.43$ Avg, Last bằng V2. Lợi thế tập trung ở bước 1--2.
\item Trần của việc chỉ sửa ``đám mây'' với backbone cố định đo được khoảng $+0.5$ Avg; muốn tăng nhiều hơn phải cải
      thiện chính bộ trích đặc trưng (LBR là một thử nghiệm theo hướng này, chưa thành công).
\end{itemize}

\section{Kết luận}
Chỉ với một bộ phân loại CA và một lần huấn luyện, HRC (mượn hình dạng lớp gần nhất, nới rộng $s=2$, $B=1024$ điểm
mỗi lớp) đạt $75.46/66.29$ trên ImageNet-A seed 1993, cao hơn V2 chỉ CA ($75.03/66.29$) và RSIAT gốc ($74.76/65.77$),
với backbone trùng khít V2. Mã: \texttt{models/hrc.py}, \texttt{models/base.py} (\texttt{\_ca\_private}); cấu hình
\texttt{exps/hrc/imageneta\_hrc\_b1024\_s1993.json}; tune ngoại tuyến \texttt{exps/hrc/tune.py}.

\begin{thebibliography}{99}
\bibitem{ssiat} Y. Tan et al. Semantically-Shifted Incremental Adapter-Tuning is A Continual ViTransformer. \emph{CVPR}, 2024.
\bibitem{slca} G. Zhang et al. SLCA: Slow Learner with Classifier Alignment for Continual Learning on a Pre-trained Model. \emph{ICCV}, 2023.
\bibitem{sdc} L. Yu et al. Semantic Drift Compensation for Class-Incremental Learning. \emph{CVPR}, 2020.
\bibitem{cosface} H. Wang et al. CosFace: Large Margin Cosine Loss for Deep Face Recognition. \emph{CVPR}, 2018.
\bibitem{freelunch} D. Goswami et al. FeCAM / Calibrating Higher-Order Statistics for Few-Shot Class-Incremental Learning with Pre-trained Vision Transformers. 2023--2024.
\bibitem{devroye} L. Devroye, L. Györfi, G. Lugosi. \emph{A Probabilistic Theory of Pattern Recognition}. Springer, 1996.
\bibitem{nc} V. Papyan, X. Y. Han, D. Donoho. Prevalence of Neural Collapse during the Terminal Phase of Deep Learning Training. \emph{PNAS}, 2020.
\bibitem{efron} B. Efron. The Estimation of Prediction Error: Covariance Penalties and Cross-Validation. \emph{JASA}, 2004.
\end{thebibliography}
\end{document}
